% Library CV materials_fr — transcription of writing-samples/cvs/Moumen_Nermine_CV_Matériaux.pdf (CV05, 2026-11-06)
\documentclass[10pt]{article}
\usepackage[letter]{nermine-cv}
\begin{document}
\cvheader{Moumen Nermine}{Étudiante en dernière année du cycle ingénieur en Chimie Industrielle}{nermine.moumen@insat.ucar.tn}{+21627040414}{linkedin.com/in/nermine-moumen}

\section{FORMATION}
\entry{Institut National des Sciences Appliquées et de Technologie (INSAT)}{2022–Présent}
\summary{Étudiante en Chimie Industrielle, spécialité Procédés et Environnement, avec une formation solide en catalyse chimique, cinétique électrochimique, science des matériaux, complétée par des bases en chimie analytique, chimie organique et HSE en laboratoire. Intéressée par la chimie verte et la valorisation de la biomasse. À la recherche d'un PFE de 6 mois à partir de février 2027.}

\section{EXPÉRIENCE PROFESSIONNELLE}
\entry{SEGULA Technologies - Stage en Modélisation Électrochimique des batteries}{06/2026–08/2026}
\begin{bul}
\item Modélisation d'une cellule \textbf{Li-ion} (NMC811/Graphite-SiOx) via le Single Particle Model (SPM) implémenté sous MATLAB.
\item Optimisation du maillage et des valeurs des coefficients de diffusion, réduisant l'erreur \textbf{RMSE de 202,13 mV à 40,5 mV}.
\item Validation des résultats par une approche itérative et une analyse critique avec les données expérimentales et COMSOL.
\end{bul}
\entry{Control Union - Stage en Contrôle Qualité}{08/2025–09/2025}
\begin{bul}
\item Analyse de lubrifiants (viscosité et métaux lourds par \textbf{ICP-OES}).
\item Préparation et interprétation de rapports d'analyse pour vérifier la conformité aux normes \textbf{ASTM D5185} et \textbf{ISO 17025}.
\end{bul}
\entry{ETAP - Stage en Caractérisation Analytique d'hydrocarbures et en Traitement des Eaux}{07/2024–08/2024}
\begin{bul}
\item Formation aux techniques de géochimie organique (GC, GC-MS, pyrolyse Rock-Eval) pour la caractérisation d'hydrocarbures.
\item Analyses physico-chimiques des eaux (pH, salinité, bicarbonates, sulfates, chlorures) par chromatographie ionique et ICP-AES.
\end{bul}

\section{PROJETS ACADÉMIQUES}
\project{Polymères à Impression Moléculaire (MIPs) en Applications Pharmaceutiques}
\begin{bul}
\item Synthèse et caractérisation d'un MIP sélectif au sulpiride par polymérisation radicalaire, avec extraction Soxhlet, suivi par spectrophotométrie UV-Vis et validation par cinétique de libération contrôlée (240 min).
\item Synthèse d'un MIP core-shell MOF/polymère pour la délivrance ciblée de doxorubicine, avec \textbf{85 \% de taux d'encapsulation}.
\item Modélisation des isothermes d'adsorption (bi-Langmuir, Freundlich, $R^2 = 0{,}9877$) pour relier les propriétés des matériaux, leurs sites de liaison et leurs performances.
\end{bul}
\project{Étude Bibliographique : Synthèse et Caractérisation de Matériaux Supraconducteurs}
\begin{bul}
\item Analyse des propriétés fondamentales des supraconducteurs et étude des voies d'élaboration : synthèse céramique à l'état solide, croissance cristalline (flux, Bridgman, Czochralski), et dépôt de couches minces (ALD).
\item Étude des méthodes de caractérisation : RAMAN, FTIR, XRD, MEB et MET.
\end{bul}
\project{Analyse techno-environnementale de la production d'éthylène : vapocraquage vs déshydratation du bioéthanol}
\begin{bul}
\item Simulation sous \textbf{Aspen Plus V14} des deux voies, dont la déshydratation du bioéthanol sur \textbf{catalyseur HZSM-5}.
\item \textbf{ACV} cradle-to-gate (OpenLCA, ISO 14040/14044) sur 5 catégories d'impact ; identification de l'origine de la biomasse et du changement d'affectation des sols indirect comme principales incertitudes.
\end{bul}
\project{Prédiction des Équilibres Liquide–Liquide pour Systèmes Ternaires d'Extraction}
\begin{bul}
\item Formulation du système thermodynamique LLE et résolution via Rachford-Rice et continuation par longueur d'arc.
\item Implémentation du workflow en Python, MATLAB et Mathematica (spline/fsolve, SciPy/mpltern, Interpolation/FindRoot).
\end{bul}

\section{COMPÉTENCES TECHNIQUES}
\begin{bul}
\item \textbf{Techniques de Laboratoire :} Chromatographie (HPLC, GC, CCM), Spectrométrie (MS, ICP-OES), Spectroscopie (RMN, FTIR, UV-Vis), Diffraction (XRD), Microscopie (MEB, MET), Électrochimie (Voltamétrie cyclique, polarographie).
\item \textbf{Simulation :} Aspen Plus et HYSYS, COMSOL Multiphysics\sep\textbf{Programmation :} \textit{MATLAB, Mathematica, Python}
\item \textbf{Données \& Bureautique :} Pack Microsoft Office, OpenLCA\sep\textbf{Normes :} ISO 9001, ISO 14001, ISO 45001, ISO 17025, ISO 14040
\end{bul}

\section{ACTIVITÉS EXTRASCOLAIRES}
\textbf{IEEE INSAT} — Chef d'organisation de TRSYP 3.0 ; Responsable des activités sociales – MENASYP ; Présidente de l'EMBS.\\
\textbf{THEATRO INSAT} — Membre de l'équipe de danse.\par

\section{LANGUES ET CERTIFICATS}
\twocol{Langues}{Anglais (Avancé), Français (Bilingue), Arabe (Natif).}{Certificats}{Aspen HYSYS – MATLAB}
\end{document}
